%%%%%%%%%%%%%%%%%%%%%%%%%%%%%%%%%%%%%%%
% Cover Letter - STI-Consulting, Werkstudent AI & Automatisierung
%%%%%%%%%%%%%%%%%%%%%%%%%%%%%%%%%%%%%%%

\documentclass[]{cover}
\usepackage{fancyhdr}

\pagestyle{fancy}
\fancyhf{}

\rfoot{Page \thepage \hspace{0pt}}
\thispagestyle{empty}
\renewcommand{\headrulewidth}{0pt}
\begin{document}

%%%%%%%%%%%%%%%%%%%%%%%%%%%%%%%%%%%%%%
%     TITLE NAME
%%%%%%%%%%%%%%%%%%%%%%%%%%%%%%%%%%%%%%
\namesection{}{\Huge{Ana Lorena Ortiz Loyola}}{  \href{mailto:a01750283@tec.mx}{a01750283@tec.mx} | +49 176 62145797 |  \urlstyle{same}\href{https://linkedin.com/in/lorenasolana}{LinkedIn}
}

%%%%%%%%%%%%%%%%%%%%%%%%%%%%%%%%%%%%%%
%     MAIN COVER LETTER CONTENT
%%%%%%%%%%%%%%%%%%%%%%%%%%%%%%%%%%%%%%

\currentdate{\today}
\lettercontent{Sehr geehrte Damen und Herren,}

\lettercontent{für die Werkstudentenstelle AI \& Automatisierung im Workstream AI-First bringe ich genau das mit, was Sie im Profil beschreiben: praktische, tägliche Arbeit mit modernen AI-Tools und ein Verständnis von Prompting als Engineering-Disziplin, nicht als einmaligen Versuch.}

\lettercontent{Besonders reizt mich, dass STI-Consulting mit dem Programm Phoenix die produktive Nutzung von KI intern selbst vorantreibt, nicht nur als Beratungsthema für Kunden. Genau dieses Zusammenspiel aus konkreter Automatisierungsarbeit und Wissensvermittlung an Kolleginnen und Kollegen motiviert mich.}

\lettercontent{Konkret bringe ich mit:}

{\raggedright\fontspec[Path = OpenFonts/fonts/raleway/]{Raleway-Medium}\fontsize{11pt}{13pt}\selectfont
\begin{itemize}
    \item \textbf{Tägliche Praxis mit AI-Tools}: In einem eigenen Projekt konzipierte und dokumentierte ich Prompts und Arbeitsabläufe mit Claude Code als KI-gestütztem Programmierassistenten, von der ersten Idee bis zur fertigen, nachvollziehbaren Datenpipeline.
    \item \textbf{Automatisierung produktiver Prozesse}: Bei der Deutschen Bundesbank (AnaCredit) entwickelte ich in Python automatisierte Tests und Fehlererkennung für regulatorische Daten, inklusive automatisierter Jira-Ticket-Erstellung und Benachrichtigungen an Fachbereiche.
    \item \textbf{Strukturierte Dokumentation}: Ich dokumentiere Arbeitsschritte und Ergebnisse gerne so, dass sie auch von Kolleginnen und Kollegen ohne technischen Hintergrund nachvollzogen werden können, eine Grundlage für die Playbooks und Schulungsunterlagen, die Sie beschreiben.
\end{itemize}\par}
\vspace{6pt}

\lettercontent{Eine Sache möchte ich offen ansprechen: Mein Bachelorabschluss ist für Januar 2027 vorgesehen, das entspricht ab jetzt rund sechs Monaten und nicht den genannten zwölf Monaten Mindeststudienzeit. Da die fachliche Passung aus meiner Sicht sehr stark ist, wollte ich mich dennoch vorstellen und bespreche den zeitlichen Rahmen gerne offen in einem Gespräch.}

\lettercontent{Ich arbeite strukturiert und eigenständig, mit Freude an der Vermittlung von Wissen, und freue mich darauf, mehr über die Rolle zu erfahren.}

\begin{flushright}
\closing{Mit freundlichen Grüßen,}

\signature{Ana Lorena Ortiz Loyola}
\end{flushright}
\end{document}
